\arrayrulecolor{Periwinkle} 

\renewcommand{\arraystretch}{1.75} % Optional: adjust row height for better readability
\setlength{\tabcolsep}{8pt} % Optional: adjust column separation for better spacing
\setlength{\arrayrulewidth}{0.75pt}


\section{SU2: Complexity}

\subsection{Introduction}

\begin{itemize}
    \item compare different algorithm complexity in order to develop more efficient programs
    \item two important aspects when comparing algorithms:
        \begin{enumerate}
            \item the \textcolor{Periwinkle}{memory} of memory used by the algorithm 
                \begin{itemize}
                    \item[] memory is determined by the chosen variables and can generally be determined by looking at the code
                \end{itemize}
            \item the \textcolor{Periwinkle}{running time} of the algorithm
                \begin{itemize}
                    \item[] running time depends on the algorithm and the number of data items processed by the algorithm 
                \end{itemize}
        \end{enumerate}
    \item the focus is on the \textbf{running time} of an algorithm - we are using algorithm analysis rather than machine runtime
\end{itemize}

\subsection{Overall Perspective}

\begin{itemize}
    \item if the running time of a programs is represented in terms of a mathematical formula, you can compare the running times of programs to each other 
    \item the table below has the mathematical formulas for the running times of \textbf{three} different programs performing the \textbf{same} task - all three algorithms create the exact same output 
    \item the graph shows the running time $T(n)$ in terms of the number of items$(n)$ handled by the algorithm
\end{itemize}

\noindent
\begin{minipage}{0.6\textwidth}
    \includegraphics[width=0.9\textwidth]{../Summaries/Images/2_1.pdf}
\end{minipage}
\begin{minipage}{0.4\textwidth}
    \begin{center}
        \renewcommand{\arraystretch}{2.3}
        \begin{tabular}{|>{\columncolor{Periwinkle!10}}l | >{\columncolor{Periwinkle!10}}c|}
            \hline
            \multicolumn{1}{|>{\columncolor{Periwinkle!30}}c|}{\textbf{Program}} & \multicolumn{1}{>{\columncolor{Periwinkle!30}}c|}{\textbf{$T(n)$}}\\
            \hline
            \multicolumn{1}{| >{\columncolor{Periwinkle!30}}l|}{Program 1} & $\displaystyle \left(\frac{11}{2}n^2 + \frac{47}{2}n + 24\right)$ \\
            \multicolumn{1}{|>{\columncolor{Periwinkle!30}}l|}{{Program 2}} & $\displaystyle 13n + 22$\\
            \multicolumn{1}{|>{\columncolor{Periwinkle!30}} c|}{Program 3} & $\displaystyle 19\left[\left(\log_2^(n + 1) \right) + 1\right] + 18$\\
            \hline
        \end{tabular}
    \end{center}
\end{minipage}

\begin{itemize}
    \item three methods to analyse the running time of algorithms:
        \begin{enumerate}
            \item the \textcolor{Periwinkle}{detailed method} using tau-notation
            \item the \textcolor{Periwinkle}{simplified model} where each tau is assigned the value equal to 1
            \item \textcolor{Periwinkle}{asymptotic analysis} using Big Oh notation in which we write $O(n)$
        \end{enumerate}
\end{itemize}

\subsubsection{Detailed Model}

\begin{itemize}
    \item \textbf{basic assumption:} the time it takes to \emph{fetch} or \emph{store} a variable in memory is always the same
\end{itemize}

\begin{center}
    \begin{minipage}{0.20\textwidth}
        \begin{tightcode}
x = y; // or 
x = 1
        \end{tightcode}
    \end{minipage}
\end{center}


\begin{itemize}
    \item[]{}
        \begin{itemize}
            \item the analysis includes the time to \textbf{fetch}
            \item \mintinline{java}{y} is the first instance; or
            \item the constant \mintinline{java}{1} in the second instance
            \item \textbf{plus} the time to \textbf{store} this value in the variable/value \mintinline{java}{x}
            \item we \textbf{do not fetch} the variable \mintinline{java}{x} as well
            \item the analysis of this line: $\displaystyle \tau_{\text{fetch}} + \tau_{\text{store}}$
        \end{itemize}
\end{itemize}

\newpage

\begin{center}
    \begin{tabular}{|
>{\columncolor{Periwinkle!10}}l| 
>{\columncolor{Periwinkle!10}}l|}
        \hline
        \multicolumn{1}{|>{\columncolor{Periwinkle!30}}c|}{\textbf{notation}} & \multicolumn{1}{>{\columncolor{Periwinkle!30}}c|}{\textbf{meaning}}\\
        $\tau_{\text{fetch}}$ & time to fetch the operand from memory\\
        $\tau_{\text{store}}$ & store the result in memory \\
        $\tau_{(+,-,\div,\times)}$ & time for basic arithmetic operations\\
        $\tau_{\leq}$ & time for all comparisons - we use the same symbol for $<$, $>$, ==, $\leq$, $\geq$\\
        $\tau_{\text{call}}$ & time to pass a parameter is constant and equal to $\tau_{\text{store}}$\\
        $\tau_{\text{return}}$ & time to call a method and the time it takes to return from a method\\
        $\tau_{[\cdot]}$ & time for the address calculation based on the index of an array\\
        \hline
    \end{tabular}
\end{center}

\subsubsection{Examples}

\paragraph*{Example 1}
\noindent
\begin{minipage}{0.05\textwidth}
    \phantom{text}
\end{minipage}
\begin{minipage}{0.2\textwidth}
    \begin{tightcode}
y = y + 1;
    \end{tightcode}
\end{minipage}
\begin{minipage}{0.24\textwidth}
    \begin{mybox}[Periwinkle]{}
        $2\tau_{\text{fetch}} + \tau_{+}+  \tau_{\text{store}}$
    \end{mybox}
\end{minipage}

\noindent
\begin{minipage}{0.05\textwidth}
    \phantom{text}
\end{minipage}
\begin{minipage}{0.2\textwidth}
    \begin{tightcode}
y = f(x);
    \end{tightcode}
\end{minipage}
\begin{minipage}{0.3\textwidth}
    \begin{mybox}[Periwinkle]{}
        $\tau_{\text{fetch}} + 2\tau_{\text{store}} + \tau_{\text{call}} + T_{f(x)}$
    \end{mybox}
\end{minipage}
\begin{minipage}{0.4\textwidth}
    \begin{mybox}[Gray]{}
        $T_{f(x)}$: running time of method $f(x)$
    \end{mybox}
\end{minipage}

\noindent
\begin{minipage}{0.05\textwidth}
    \phantom{text}
\end{minipage}
\begin{minipage}{0.2\textwidth}
    \begin{tightcode}
y = a[i];
    \end{tightcode}
\end{minipage}
\begin{minipage}{0.22\textwidth}
    \begin{mybox}[Periwinkle]{}
        $3\tau_{\text{fetch}} + \tau_{[\cdot]} \tau_{\text{store}}$
    \end{mybox}
\end{minipage}
\begin{minipage}{0.25\textwidth}
    \begin{mybox}[Gray]{}
        3 fetches: \mintinline{text}{i}; \mintinline{text}{a}; \mintinline{text}{a[i]}
    \end{mybox}
\end{minipage}

\paragraph*{Example 2}

\begin{center}
    \begin{minipage}{0.33\textwidth}
        \begin{tightcode}
for(int i = 1; i <= n; ++i)
        \end{tightcode}
    \end{minipage}
\end{center}

\noindent
\begin{minipage}{0.15\textwidth}
    \begin{tightcode}
i = 1;
    \end{tightcode}
\end{minipage}
\begin{minipage}{0.5\textwidth}
    only once, at the beginning of the loop\\
    $\tau_{\text{fetch}} + \tau_{\text{store}}$
\end{minipage}

\vspace{0.35cm}

\noindent
\begin{minipage}{0.15\textwidth}
    \begin{tightcode}
i <= n;
    \end{tightcode}
\end{minipage}
\begin{minipage}{0.75\textwidth}
    $2\tau_{\text{fetch}} + \tau_{<}$ - how many times will the line execute?\\
    it will start at \mintinline{java}{i = 1} and for the final time for \mintinline{java}{(i = n + 1)} $\rightarrow n + 1$ times\\
    this part is always executed one more time for the false condition
\end{minipage}

\vspace{0.35cm}

\noindent
\begin{minipage}{0.15\textwidth}
    \begin{tightcode}
++i;
    \end{tightcode}
\end{minipage}
\begin{minipage}{0.75\textwidth} 
    rewrite as \mintinline{java}{i = i + 1} $\rightarrow 2\tau_{\text{fetch}} + \tau_{+} + \tau_{\text{store}}$\\
    it will start at \mintinline{java}{i = n} and the final time for \mintinline{java}{i = n}, which is $n$ times
\end{minipage}

\begin{center}
    therefore, the total running time for this \mintinline{java}{for} loop is:\\
    $\tau_{\text{fetch}} + \tau_{\text{store}} + (n+1)(2\tau_{\text{fetch}} + \tau_{<}) + n(2\tau_{\text{fetch}} + \tau_{+}) + \tau_{\text{store}}$
\end{center}

\paragraph*{Example 3}

\begin{center}
    \begin{minipage}{0.4\textwidth}
    \begin{tightcode}
public class Example3 {
    public static int sum(int n) {
        int result = 0;
        for(int i = 1; i <=n; ++1)
            result += i;
        return result;
    }
}
    \end{tightcode}
\end{minipage}
\end{center}


\begin{center}
    \begin{tabular}{| 
        >{\columncolor{Periwinkle!10}}c | 
        >{\columncolor{Periwinkle!10}}l |
        >{\columncolor{Periwinkle!10}}l |}
        \hline
        \multicolumn{1}{|>{\columncolor{Periwinkle!30}}c|}{\textbf{statement}} & \multicolumn{1}{>{\columncolor{Periwinkle!30}}c|}{\textbf{code}} & \multicolumn{1}{>{\columncolor{Periwinkle!30}}c|}{\textbf{time}}\\
        \hline
        3 &  \mintinline{java}{result = 0} & $\tau_{\text{fetch}} + \tau_{\text{store}}$\\
        4a &   \mintinline{java}{i = 1} & $\tau_{\text{fetch}} + \tau_{\text{store}}$\\
        4b &  \mintinline{java}{i <= n} & $(n+1)(2\tau_{\text{fetch}} + \tau_{\text{store}})$\\
        4c &  \mintinline{java}{++i} & $(n)(2\tau_{\text{fetch}} + \tau_{<} + \tau_{\text{store}})$\\
        5 &  \mintinline{java}{result += i} & $(n)(2\tau_{\text{fetch}} + \tau_{<} + \tau_{\text{store}})$\\
        6 &  \mintinline{java}{return result} & $\tau_{\text{fetch}} + \tau_{\text{return}}$\\
        \hline
        \hline
        \multicolumn{3}{|>{\columncolor{Periwinkle!30}}l|}{\makecell{$(n)(6\tau_{\text{fetch}} + 2\tau_{\text{store}} + \tau_{<} + \tau_{\text{+}}) + (5\tau_{\text{fetch}} + 2\tau_{\text{store}} + \tau_{<}) + \tau_{\text{return}}$}}\\
        \hline
    \end{tabular}
\end{center} 

\paragraph*{Example 4}

\begin{center}
    \begin{minipage}{0.6\textwidth}
        \begin{tightcode}
public class Example4 {
    public static int horner(int[] a, int n, int x) {
        int result = a[n];
        for(int i = n - 1; i >= 0; --i)
            result = result * x + a[i];
        return result
    }
}
        \end{tightcode}
    \end{minipage}
\end{center}

\begin{center}
    \begin{tabular}{| 
        >{\columncolor{Periwinkle!10}}c | 
        >{\columncolor{Periwinkle!10}}l |
        >{\columncolor{Periwinkle!10}}l |}
        \hline
        \multicolumn{1}{|>{\columncolor{Periwinkle!30}}c|}{\textbf{statement}} & \multicolumn{1}{>{\columncolor{Periwinkle!30}}c|}{\textbf{code}} & \multicolumn{1}{>{\columncolor{Periwinkle!30}}c|}{\textbf{time}}\\
        \hline
        3 &  \mintinline{java}{int result = a[n]} & $3\tau_{\text{fetch}} + \tau_{[\cdot]} + \tau_{\text{store}}$\\
        4a &  \mintinline{java}{int i = n - 1} & $2\tau_{\text{fetch}} + \tau_{-} + \tau_{\text{store}}$\\
        4b &  \mintinline{java}{i >= 0} & $(n+1)(2\tau_{\text{fetch}} + \tau_{<})$\\
        4c &  \mintinline{java}{--i} & $(n)(2\tau_{\text{fetch}} + \tau_{-} + \tau_{\text{store}})$\\
        5 &  \mintinline{java}{result = result * x + a[i]} & $(n)(5\tau_{\text{fetch}} + \tau_{[\cdot]} + \tau_{+} + \tau_{\times} + \tau_{\text{store}})$\\
        6 &  \mintinline{java}{return result} & $\tau_{\text{fetch}} + \tau_{\text{store}}$\\
        \hline \hline
        \multicolumn{3}{|>{\columncolor{Periwinkle!30}}l|}{$(n)(9\tau_{\text{fetch}} + 2\tau_{\text{store}} + \tau_{<} + \tau_{[\cdot]}+ \tau_{+} + \tau_{\times} + \tau_{-}) + (7\tau_{\text{fetch}} + 2\tau_{\text{store}} + \tau_{[\cdot]} + \tau_{-} + \tau_{<} + \tau_{\text{return}})$}\\
        \hline
    \end{tabular}
\end{center} 

\subsection{More on  \text{\mintinline{java}{for}} loops}

\begin{center}
    \noindent
    \begin{minipage}{0.32\textwidth}
            \begin{tightcode}
for(int i = 0; i < n; i++)
        \end{tightcode}
    \end{minipage}
\end{center}


\begin{center}
    \renewcommand{\arraystretch}{2.5}
    \begin{tabular}{l l l}
        \textbf{\# of times part b is executed:}  & $n+1$ & one more time than \textbf{c}\\
        \textbf{\# of times part c is executed:} & $n$ & from 0 to $n - 1$ is $n$ times\\
        \textbf{notes:} if $n = 3$ & \multicolumn{2}{l}{\makecell[l]{the valid values are 0; 1; 2\\
    this is three values $\therefore n$ times}}
    \end{tabular}
\end{center}

% \renewcommand{\arraystretch}{1.75}

\begin{center}
    \noindent
    \begin{minipage}{0.28\textwidth}
        \begin{tightcode}
for(i = 1; i < n; i++)
        \end{tightcode}
    \end{minipage}
\end{center}

\newpage

\begin{center}
    \begin{tabular}{l l l}
        \textbf{\# of times part b is executed:}  & $n$ & one more time than \textbf{c}\\
        \textbf{\# of times part c is executed:} & $n$ & from 1 to $(n - 1)$ is $n$ times\\
        \textbf{notes:} if $n=3$& \multicolumn{2}{l}{\makecell[l]{the valid values are 1; 2\\this is two values $\therefore (n-1)$ times}}
    \end{tabular}
\end{center}

\begin{center}
    \begin{minipage}{0.29\textwidth}
        \begin{tightcode}
for(i = 0; i <= n; i++)
        \end{tightcode}
    \end{minipage}
\end{center}


\begin{center}
    \begin{tabular}{l l l}
        \textbf{\# of times part b is executed:}  & $n + 2$ & one more time than \textbf{c}\\
        \textbf{\# of times part c is executed:} & $n+1$ & from 0 to $n$ is $(n + 1)$ times\\
        \textbf{notes:} if $n=3$& \multicolumn{2}{l}{\makecell[l]{the valid values are 0; 1; 2; 3\\this is four values $\therefore (n+1)$ times}}
    \end{tabular}
\end{center}

\begin{center}
    \begin{minipage}{0.29\textwidth}
        \begin{tightcode}
for(i = 1; i <= n; i++)
        \end{tightcode}
    \end{minipage}
\end{center}

\begin{center}
    \begin{tabular}{l l l}
        \textbf{\# of times part b is executed:}  & $n + 1$ & one more time than \textbf{c}\\
        \textbf{\# of times part c is executed:} & $n$ & from 1 to $n$ is $n$ times\\
        \textbf{notes:} if $n=3$& \multicolumn{2}{l}{\makecell[l]{the valid values are 1; 2; 3\\this is three values $\therefore n$ times}}
    \end{tabular}
\end{center}

\begin{center}
    \begin{minipage}{0.35\textwidth}
        \begin{tightcode}
for(i = 2; i <= (n + 2); i++)
        \end{tightcode}
    \end{minipage}
\end{center}

\begin{center}
    \begin{tabular}{l l l}
        \textbf{\# of times part b is executed:}  & $n + 2$ & one more time than \textbf{c}\\
        \textbf{\# of times part c is executed:} & $n+1$ & from 1 to $n$ is $(n+1)$ times\\
        \textbf{notes:} if $n=3$& \multicolumn{2}{l}{\makecell[l]{the valid values are 2; 3; 4; 5\\this is four values $\therefore n+1$ times}}
    \end{tabular}
\end{center}


\raggedright \includegraphics[width=1\textwidth]{../Summaries/Images/2_2.pdf}

\subsection{Simplified Model}

\begin{itemize}
    \item assume all the $\tau$ values are taking the same time and that this time is equal to 1
    \item $\tau_{\text{fetch}} == \tau_{\text{store}} == \tau_{\cdots} == 1$
    \item the resulting equation can be drawn on a graph 
    \item note the difference between the fixed time (5) and the variable time (7)
    \item revisit the following \mintinline{java}{for} loop:
\end{itemize}

\begin{center}
    \begin{minipage}{0.33\textwidth}
        \begin{tightcode}
for(int i = 1; i <= n; i++)
        \end{tightcode}
    \end{minipage}
\end{center}

\begin{itemize}
    \item[] \textbf{detailed analysis:} 
\end{itemize}

\begin{center}
    $\displaystyle \tau_{\text{fetch}} + \tau_{\text{store}}+ (n+1)(2\tau_{\text{fetch}} + \tau_{<}) + n(2\tau_{\text{fetch}} + \tau_{+} + \tau_{\text{store}})$
\end{center}

\begin{itemize}
    \item[] \textbf{simplified model:}
\end{itemize}

\vspace{-0.75cm}

\begin{align*}
    & (1 + 1) + (n + 1)(2 + 1) + n(2 + 1 + 1)\\
    & = (2) + (n+1)(3) + n(4)\\
    & = 2 + 3n + 3 + 4n\\
    & = 7n + 5
\end{align*}

\subsubsection{The Double \text{\mintinline{java}{for}} Loop}

\noindent
\begin{minipage}{0.62\textwidth}
    \begin{center}
        \begin{tightcode}
public class Example1 {
    public static int geometricSeriesSum(int x, int n) {
        int sum = 0;
        for(int i = 0; i <= n; ++i) {
            int prod = 1;
            for (int j = 0; j < i; ++j) 
                prod *= x;
            sum += prod
        }
        return sum;
    }
}

        \end{tightcode}
    \end{center}
\end{minipage}
\begin{minipage}{0.35\textwidth}
    \begin{center}
        \begin{tabular}{|>{\columncolor{Periwinkle!10}} c | 
            >{\columncolor{Periwinkle!10}} l |
            >{\columncolor{Periwinkle!10}} c |}
            \hline
            \multicolumn{1}{|>{\columncolor{Periwinkle!30}}c|}{\textbf{statement}} &
            \multicolumn{1}{>{\columncolor{Periwinkle!30}}c|}{\textbf{code}} &
            \multicolumn{1}{>{\columncolor{Periwinkle!30}}c|}{\textbf{time}}\\
            \hline
            3 & \mintinline{java}{int sum = 0} & 2\\
            4a & \mintinline{java}{int i = 0} & 2\\
            4b & \mintinline{java}{i <= n} & $3(n+2)$\\
            4c & \mintinline{java}{++i} & 4(n+1)\\
            5 & \mintinline{java}{int prod = 1} & $2(n+1)$\\
            8 & \mintinline{java}{sum += prod} & 4(n+1)\\
            6 & \mintinline{java}{return result} & 2\\
            \hline
        \end{tabular}
    \end{center}
\end{minipage}

\begin{itemize}
    \item the summation's first and last values are the first and last valid values of the counter from the \textcolor{Periwinkle}{outer} \mintinline{java}{for} loop
    \item inside the brackets is the number of times the \textcolor{Periwinkle}{inner} loop executes - note that the variable \mintinline{java}{i} \textbf{stops the iteration, NOT at $n$} as in the outer loop 
        \begin{itemize}
            \item so the outer loop is from \mintinline{java}{j = 0} to \mintinline{java}{j < i}, which is $i$ times for the \textbf{c} part, but 1 more for the \textbf{b} part, which is $(i + 1)$
        \end{itemize}
\end{itemize}

\noindent
\begin{minipage}{0.62\textwidth}
    \begin{center}
        \begin{tightcode}
int sum = 0;
    for(int i = 0; i <= n; ++i) {
        int prod = 1;
        for (int j = 0; j < i; ++j) {
            prod *= x;
        }
        sum += prod
    }
return sum
        \end{tightcode}
    \end{center}
\end{minipage}
\begin{minipage}{0.35\textwidth}
    \begin{center}
        \begin{tabular}{
            |>{\columncolor{Periwinkle!10}} c | 
            >{\columncolor{Periwinkle!10}} l |
            >{\columncolor{Periwinkle!10}} c |}
            \hline
            \multicolumn{1}{|>{\columncolor{Periwinkle!30}}c|}{\textbf{statement}} &
            \multicolumn{1}{>{\columncolor{Periwinkle!30}}c|}{\textbf{code}} &
            \multicolumn{1}{>{\columncolor{Periwinkle!30}}c|}{\textbf{time}} 
            \\
            \hline
            1 & \mintinline{java}{int sum = 0} & 2\\
            2a & \mintinline{java}{int j = 0} & 2\\
            2b & \mintinline{java}{j <= n} & $3(n + 2)$\\
            2c & \mintinline{java}{++j} & $4(n + 1)$\\
            3 & \mintinline{java}{int prod = 1} & $2(n+1)$\\
            7 & \mintinline{java}{sum += prod} & $4(n + 1)$\\
            9 & \mintinline{java}{return sum} & 2\\
            \hline
        \end{tabular}
    \end{center}
\end{minipage}

\begin{itemize}
    \item the number before the summation is the analysis of the code itself - for \textbf{6b} it is $2\tau_{\text{fetch}} + \tau_{<}$, which adds up to three units
    \item remember, line \textbf{6a} is not repeated in the loop - so lines \textbf{5} and \textbf{6a}
    \item line \textbf{7} is similar to line \textbf{6c} in terms of how many times it is executed - because the the inner part of the loop is executed the same number of times as the \textbf{c} part of the statement in the \mintinline{java}{for} statement
    \item a \mintinline{java}{for} loop leads to an equation dependant on $n$
    \item a \textbf{double} \mintinline{java}{for} loop leads to an equation dependant on $n^2$
\end{itemize}























\begin{comment}



\includegraphics[width=0.9\textwidth]{../Summaries/Images/2_1.pdf}


    \mintinline{}{}



\tau_{\text{}}

>{\columncolor{Periwinkle!10}}

\end{comment}